%-----------------------------------------------------------------------------------------------------------------------------------------------%
%	The MIT License (MIT)
%
%	Copyright (c) 2021 Jitin Nair
%
%	Permission is hereby granted, free of charge, to any person obtaining a copy
%	of this software and associated documentation files (the "Software"), to deal
%	in the Software without restriction, including without limitation the rights
%	to use, copy, modify, merge, publish, distribute, sublicense, and/or sell
%	copies of the Software, and to permit persons to whom the Software is
%	furnished to do so, subject to the following conditions:
%	
%	THE SOFTWARE IS PROVIDED "AS IS", WITHOUT WARRANTY OF ANY KIND, EXPRESS OR
%	IMPLIED, INCLUDING BUT NOT LIMITED TO THE WARRANTIES OF MERCHANTABILITY,
%	FITNESS FOR A PARTICULAR PURPOSE AND NONINFRINGEMENT. IN NO EVENT SHALL THE
%	AUTHORS OR COPYRIGHT HOLDERS BE LIABLE FOR ANY CLAIM, DAMAGES OR OTHER
%	LIABILITY, WHETHER IN AN ACTION OF CONTRACT, TORT OR OTHERWISE, ARISING FROM,
%	OUT OF OR IN CONNECTION WITH THE SOFTWARE OR THE USE OR OTHER DEALINGS IN
%	THE SOFTWARE.
%	
%
%-----------------------------------------------------------------------------------------------------------------------------------------------%

%----------------------------------------------------------------------------------------
%	DOCUMENT DEFINITION
%----------------------------------------------------------------------------------------

\documentclass[a4paper,12pt]{article}

\usepackage{ctex} 

%----------------------------------------------------------------------------------------
%	FONT
%----------------------------------------------------------------------------------------

% % fontspec allows you to use TTF/OTF fonts directly
% \usepackage{fontspec}
% \defaultfontfeatures{Ligatures=TeX}

%----------------------------------------------------------------------------------------
%	PACKAGES
%----------------------------------------------------------------------------------------
\usepackage{url}
\usepackage{parskip} 	

%other packages for formatting
\RequirePackage{color}
\RequirePackage{graphicx}
\usepackage[usenames,dvipsnames]{xcolor}
\usepackage[scale=0.9]{geometry}

%tabularx environment
\usepackage{tabularx}

%for lists within experience section
\usepackage{enumitem}

% centered version of 'X' col. type
\newcolumntype{C}{>{\centering\arraybackslash}X} 

%to prevent spillover of tabular into next pages
\usepackage{supertabular}
\usepackage{tabularx}
\newlength{\fullcollw}
\setlength{\fullcollw}{0.47\textwidth}

%custom \section
\usepackage{titlesec}				
\usepackage{multicol}
\usepackage{multirow}

%CV Sections inspired by: 
%http://stefano.italians.nl/archives/26
\titleformat{\section}{\Large\scshape\raggedright}{}{0em}{}[\titlerule]
\titlespacing{\section}{0pt}{10pt}{10pt}

%for publications
\usepackage[style=authoryear,sorting=ynt, maxbibnames=2]{biblatex}

%Setup hyperref package, and colours for links
\usepackage[unicode, draft=false]{hyperref}
\definecolor{linkcolour}{rgb}{0,0.2,0.6}
\hypersetup{colorlinks,breaklinks,urlcolor=linkcolour,linkcolor=linkcolour}
\addbibresource{citations.bib}
\setlength\bibitemsep{1em}

%for social icons
\usepackage{fontawesome5}

%debug page outer frames
%\usepackage{showframe}


% job listing environments
\newenvironment{jobshort}[2]
    {
    \begin{tabularx}{\linewidth}{@{}l X r@{}}
    \textbf{#1} & \hfill &  #2 \\[3.75pt]
    \end{tabularx}
    }
    {
    }

\newenvironment{joblong}[2]
    {
    \begin{tabularx}{\linewidth}{@{}l X r@{}}
    \textbf{#1} & \hfill &  #2 \\[3.75pt]
    \end{tabularx}
    \begin{minipage}[t]{\linewidth}
    \begin{itemize}[nosep,after=\strut, leftmargin=1em, itemsep=3pt,label=--]
    }
    {
    \end{itemize}
    \end{minipage}    
    }

%----------------------------------------------------------------------------------------
%	BEGIN DOCUMENT
%----------------------------------------------------------------------------------------

\begin{document}

\pagestyle{empty} 

%----------------------------------------------------------------------------------------
%	TITLE
%----------------------------------------------------------------------------------------

\begin{tabularx}{\linewidth}{@{} C @{}}
\Huge{左天佑} \\[7.5pt]
\href{mailto:fztym@sjtu.edu.cn}{\raisebox{-0.05\height}\faEnvelope \ fztym@sjtu.edu.cn} \ $|$ \ 
\href{tel:+8613162909941}{\raisebox{-0.05\height}\faMobile \ +86.131.6290.9941} \\
\end{tabularx}

%----------------------------------------------------------------------------------------
%	EDUCATION
%----------------------------------------------------------------------------------------

\section{教育背景}
\begin{tabularx}{\linewidth}{@{}l X@{}}	
\textbf{上海交通大学} & 中国上海 \\
\textit{电子计算机工程，计算机科学辅修} & 2023 年 9 月 - 至今 \\  
\normalsize (GPA: 3.62/4.0) & (排名：51/215)\\
\end{tabularx}

%----------------------------------------------------------------------------------------
%	Courses
%----------------------------------------------------------------------------------------

\section{主修课程}
\begin{tabularx}{\linewidth}{@{}l X@{}}	
ENGR1510J & 计算机导论（进阶） \hfill \normalsize (4.0/4.0)\\
MATH1860J & 荣誉数学 II \hfill \normalsize (4.0/4.0)\\
ECE2810J & 数据结构与算法 \hfill \normalsize (4.0/4.0)\\
ECE4770J & 算法导论 \hfill \normalsize (4.0/4.0)\\
ECE4720J & 大数据方法与工具 \hfill \normalsize (4.0/4.0)\\
% ECE3480J & 量子信息科学导论 \hfill \normalsize (4.0/4.0) \\
ECE4821J & 操作系统导论 \hfill \normalsize (4.0/4.0)
\end{tabularx}

%----------------------------------------------------------------------------------------
%	Prizes & Honors
%----------------------------------------------------------------------------------------

\section{奖项和荣誉}
\begin{tabularx}{\linewidth}{@{}l X@{}}
\textbf{一等奖}, & NOIP 2020 \\
\textbf{银奖}, & UPC 2024
\end{tabularx}


%----------------------------------------------------------------------------------------
% EXPERIENCE SECTIONS
%----------------------------------------------------------------------------------------

\section{工作经历}

\begin{jobshort}{摩尔线程 (Moore Threads) | GPU 计算与 AI 软件实习生}{2026 年 1 月 - 2026 年 5 月}
参与 muBLAS 高性能 GPU 数学库的开发。
实现并优化用于线性代数例程（包括 LU 分解和批处理操作）的 CUDA 内核。
分析共享内存使用情况及归约策略，以提升数值稳定性和性能。
\end{jobshort}

\begin{jobshort}{ECE4770J 算法导论 | 助教}{2025 年 9 月 - 2025 年 12 月}
开展关于 OCaml 算法设计与实现的实验课，重点讲解尾递归优化和延续传递风格 (CPS)。
维护并扩展在线判题系统 (OJ)，强化测试用例，并调试评分脚本。
\end{jobshort}

\begin{jobshort}{ECE2810J 数据结构与算法 | 助教}{2025 年 9 月 - 2025 年 12 月}
主持习题课，涵盖高级数据结构与复杂度分析。
设计受信息竞赛启发的补充习题集，并维护课程 OJ 系统。
\end{jobshort}

% \begin{jobshort}{ENGR1000J 工程导论 (SE) | 助教}{2025 年 5 月 - 2025 年 8 月}
%主持实验室课程，涵盖 Git 工作流、函数式编程、敏捷开发实践及基础游戏引擎架构。
% 维护课程仓库，并针对学生代码质量和团队协作提供结构化反馈。
% \end{jobshort}

%Projects
\section{项目经历}

\begin{tabularx}{\linewidth}{ @{}l r@{} }

\textbf{muBLAS} \\[3.75pt]
\multicolumn{2}{@{}X@{}}{muBLAS 是摩尔线程（Moore Threads）自主研发的 GPU 基础线性代数算子库（对标 NVIDIA cuBLAS）。期间参与了多项核心矩阵运算 API 的设计与开发，包括批处理矩阵 LU 分解 \mintinline{cpp}{getrfBatched}、批处理线性方程组求解 \mintinline{cpp}{getrsBatched}、以及批处理矩阵求逆 \mintinline{cpp}{getriBatched} 和 \mintinline{cpp}{matinvBatched}。}  \\

\textbf{LemonDB 现代 C++ 多线程内存数据库} & \hfill \href{https://github.com/zty2004/ECE4821Jp2}{项目链接} \\[3.75pt]
\multicolumn{2}{@{}X@{}}{基于 \textbf{C++20} 开发高并发关系型数据库。核心主导高性能拉取式（Pull-Based）线程池，通过 Worker 线程批量拉取任务（Batch Fetch）将全局锁竞争降低 16 倍，并利用本地队列减少同步开销；并发控制采用表级 \textbf{std::shared\_mutex} 配合 RAII 守卫确保异常安全。设计自适应执行路径，小负载自动降级为单线程以消除线程池开销。构建依赖感知调度系统，利用 \textbf{std::variant} 多级队列与全局优先级堆，实现复杂查询间资源依赖的自动解析与透明拓扑调度。} \\

\textbf{Mum Shell} & \hfill \href{https://focs.ji.sjtu.edu.cn/git/ece482/ZuoTianyou523370910102-p1}{项目链接} \\[3.75pt]
\multicolumn{2}{@{}X@{}}{使用 C 语言构建了一个类 Unix Shell，支持重定向、管道、后台作业及信号处理。利用 fork, exec 和 waitpid 实现了进程控制。}  \\

\textbf{MSD分布式音乐推荐系统} & \hfill \href{https://github.com/zty2004/ECE4721J}{项目链接} \\[3.75pt]
\multicolumn{2}{@{}X@{}}{基于百万歌曲数据集构建端到端大数据管道；将288GB原始数据转为884MB Avro格式（存储效率提升99.7\%），使PySpark处理耗时从5.5小时缩短至1.3小时。基于FAISS(HNSW/LSH)实现亚秒级高性能推荐算法，并利用Spark MLlib与XGBoost训练预测模型，在$\pm$10年误差内达到80\%的准确率。} \\

% \textbf{Matrix Deception} & \hfill \href{https://focs.ji.sjtu.edu.cn/silverfocs/demo/2024/p2team16/}{游戏链接} \\[3.75pt]
% \multicolumn{2}{@{}X@{}}{设计并实现了核心游戏逻辑与状态管理，采用了模块化架构以保持代码的可维护性。}  \\

% \textbf{Uno 游戏} & \hfill \href{https://focs.ji.sjtu.edu.cn/git/engr151-23fa/ZuoTianyou523370910102-p2}{游戏链接} \\[3.75pt]
% \multicolumn{2}{@{}X@{}}{独立使用 C 语言实现了完整的游戏逻辑，采用了结构化的程序设计方法。}  \\

\end{tabularx}

%----------------------------------------------------------------------------------------
%	SKILLS
%----------------------------------------------------------------------------------------

\section{技能专长}
\begin{tabularx}{\linewidth}{@{}l X@{}}
编程语言 &  \normalsize{C, C++, CUDA, OCaml, SQL, Java, Scala, Shell, Python} \\
语言能力 &  \normalsize{普通话，英语} \\
技术技能 & \normalsize{Git, SSH, CUDA 内核开发，Linux, Docker, Hadoop, Spark} \\
\end{tabularx}

\vfill
\center{\footnotesize 最后更新：\today}

\end{document}
